\section{Theoretical Investigations}
\label{sec:theoretical_analysis}

This section uses ten derived tables (Appendix~\ref{app:app2}) to examine specific questions about representation geometry, mean removal, spectral filtering, and pooling.

\subsection{Base-to-Dedicated Gap Closure Dynamics}
Table~\ref{tab:gap_closure_ir} and Table~\ref{tab:gap_closure_sts} measure how much of the gap between raw mean-pooled base models and the dedicated encoders each method closes without retraining.

Soft-ZCA closes a larger share of the gap on every STS dataset than on any IR dataset. Within IR, it closes the most on ArguAna, where the gap is smallest, and the least on FiQA2018, where the gap is largest. The best deflation method (among R1, R2, and ABTT) closes less than Soft-ZCA everywhere. It closes close to a third of the gap on STSBenchmark and SICK-R, but none on SCIDOCS and STS22.v2, where every deflation method stays below the raw baseline.

\subsection{SpecTemp Dynamics and Spectral Decay Properties}
Table~\ref{tab:spectemp_dynamics} lists the spectral properties of \texttt{google/embeddinggemma-300m} ($d=768$) on each dataset and the derived tempering exponents $\gamma(k)$.

On every dataset, $\gamma(k)$ falls as $k$ grows. It lies between about 0.4 and 0.9 at $k=64$ and is close to zero at $k=512$, where SpecTemp approaches unweighted PCA ($\gamma = 0$). The knee rank found by Kneedle lies between 23 and 55, and datasets with a later knee receive a larger $\gamma$ at $k=64$. STS22.v2 is an outlier: its leading signal-to-noise ratio is several orders of magnitude higher than on the other datasets, and its $\gamma$ at $k=512$ is zero.

\subsection{Intrinsic Space Geometry and Downstream Task Correlation}
Table~\ref{tab:geometry_task_correlation} correlates each geometry metric with task scores, pooling all models, methods, and datasets within each track.

On both tracks, higher MEV goes with lower scores and higher NID with higher scores, both moderately. IsoScore has only a weak linear relation to scores but a moderate rank relation. Centroid norm and average cosine have the weakest linear relations, and neither is significant on IR. These correlations come from a pooled set and do not show that raising isotropy raises scores: Soft-ZCA gives the dedicated encoders their highest IsoScore while lowering their scores (Section~\ref{sec:geometry_results}).

\subsection{R2 vs R1 Head-to-Head Error Cancellation}
Table~\ref{tab:r2_vs_r1_head_to_head} compares mean-subspace projection deflation (R2) with direct mean subtraction (R1). It tests the hypothesis of \citet{ren2025correcting} that R2 outperforms R1 by canceling first-order parallel sample-mean estimation noise.

R2 scores above R1 in most comparisons. On STS it falls behind only on STS22.v2, for two models. The differences are very small, mostly below 0.003 nDCG@10 on IR and below 1 Spearman point on every STS comparison, so on these benchmarks the two methods are nearly interchangeable.

\subsection{All-But-The-Top (ABTT) Monotonicity Progression}
Table~\ref{tab:abtt_monotonicity} tests whether removing successive principal components ($D \in \{1, 2, 3\}$) improves scores step by step.

A steady improvement with each removed component appears only for base models, mainly on STSBenchmark and with last-token pooling, and never for dedicated encoders. With last-token pooling, ABTT-3 beats the raw baseline for both base models on every dataset. With mean pooling, it falls below the baseline on SCIDOCS for both models, on ArguAna for Qwen3.5, and on STS22.v2 for Gemma-3. For dedicated encoders, ABTT-3 is below the baseline in every case except EmbeddingGemma on SCIDOCS.

\subsection{Dual-Regime Compression Divergence}
Table~\ref{tab:compression_divergence} contrasts SpecTemp with whitening at $k=512$ and with PCA at $k=64$.

For dedicated encoders, SpecTemp beats whitening at $k=512$ in every case except EmbeddingGemma on STS22.v2. At $k=64$ its comparison with PCA is mixed: SpecTemp is ahead for both encoders on SCIDOCS, STS22.v2, and STSBenchmark, and behind for both on ArguAna. For base models the picture reverses at $k=512$, where whitening beats SpecTemp in every case. At $k=64$, SpecTemp beats PCA for base models in every STS case and in most IR cases. SpecTemp therefore avoids the large-$k$ cost of whitening only for dedicated encoders.

\subsection{Scientific Negative Controls Validation}
Table~\ref{tab:negative_controls_validation} compares three methods with matching controls: R1 with mean centering without renormalization (\texttt{mc}), R2 with deflation of a random direction (\texttt{rand}), and PCA with random truncation at $k=128$.

Renormalization matters. R1 beats \texttt{mc} in every STS case and in most IR cases; the exceptions are EmbeddingGemma on ArguAna and last-token pooling on SCIDOCS for both base models. The largest gap is for Gemma-3 with mean pooling on STS22.v2.

R2 does not reliably beat the random direction, which itself scores almost exactly like the baseline. R2 is ahead on STSBenchmark for every model, on SICK-R for the base models, and on FiQA2018 for most models. It is behind for most models on ArguAna, SCIDOCS, and STS22.v2. Removing the mean direction is therefore not consistently better than removing a random one.

PCA beats random truncation for the dedicated encoders on every IR dataset. Random truncation is ahead on STS22.v2 for every model and on SCIDOCS for every base model.

\subsection{Sequence Length and Context Scale Impact}
Table~\ref{tab:context_length_impact} lists, for each benchmark, the last-token pooling penalty and the Soft-ZCA gain for mean-pooled base models.

The penalty is small, or even favors last-token pooling, on the two short-sentence datasets SICK-R and STSBenchmark, which average under 10 words per sentence. It is large on the three IR datasets and on STS22.v2, whose texts are much longer (Section~\ref{sec:datasets}). Among these longer-text datasets the relative penalty does not grow steadily with text length, and it is largest on SCIDOCS. The Soft-ZCA gain for mean pooling is positive on every benchmark.

\subsection{Post-Hoc Spectral Compression vs Native MRL}
Table~\ref{tab:spectemp_vs_mrl} compares SpecTemp with Matryoshka Representation Learning (MRL) prefix slicing on the dedicated encoders.

SpecTemp matches or beats prefix slicing more often as $k$ shrinks, but the outcome depends mainly on the task. On IR, SpecTemp matches or beats prefix slicing in nearly every case. The exceptions are Qwen3-Embedding on ArguAna at $k=256$ and $k=512$ and on FiQA2018 at $k=512$. On STS, prefix slicing is ahead in most cases, including Qwen3-Embedding on SICK-R at every $k$ and both encoders on STS22.v2 at every $k$ except 64. On STSBenchmark the encoders disagree: SpecTemp is ahead at every $k$ for Qwen3-Embedding but behind at every $k$ except 512 for EmbeddingGemma. SpecTemp is fitted on the evaluation corpus itself, while prefix slicing needs no data.

\subsection{Tri-Paradigm Degeneracy Benchmark}
Table~\ref{tab:tri_paradigm_benchmark} compares standardization (coordinate), R2 (subspace), and Soft-ZCA (spectral) on the base models.

Soft-ZCA has the highest nDCG@10 or Spearman $\rho$ of the three for every base model, pooling mode, and dataset. Standardization is usually ahead of R2, and most exceptions are on FiQA2018. On STS22.v2 with mean pooling, both standardization and R2 fall below the raw baseline, so only Soft-ZCA helps. The comparison covers only these three methods and the two primary metrics. On SICK-R, ABTT gives base models a higher Pearson $r$ than Soft-ZCA (Section~\ref{sec:similarity_results}).
